\chapter{Named Shape Groups and Rank-Agnostic Congruence}
\label{ch:shapes}

\chapterorient{The previous chapter introduced tensors and promised a precise
account of their \emph{shapes}. This is that account. We will build a small
notation that lets a \emph{single} function signature describe an operation on
vectors, matrices, and higher tensors all at once---a property we call
\emph{rank-agnostic congruence}---and we will see why the obvious first
attempt, writing \code{Tensor[N]} for ``the same shape as the other operand,''
silently throws that property away. By the end you will be able to read and
write the shape contracts that label every function in
\cref{part:framework}: named shape groups for congruence, the range-first
operator-shape form, the fixed-meaning axes of the matrix-free kernel, and the
one typing rule---scalar promotion---that lets a real number quietly join a
complex computation.}

\section{The problem with \texorpdfstring{\code{Tensor[N]}}{Tensor[N]}}
\label{sec:sh-problem}

In \cref{ch:tensors} we wrote a tensor's shape as a bracketed list of axis
extents: a length-$N$ vector is \lstinline[style=l4style]|Tensor[N]|, an
$m\times n$ matrix is \lstinline[style=l4style]|Tensor[m, n]|, the
distribution field of \cref{ch:bigidea} is
\lstinline[style=l4style]|Tensor[H, W, VY=3, VX=3]|. Each entry in the bracket
is one axis; the \emph{number} of entries is the tensor's \emph{rank}. A
vector has rank $1$, a matrix rank $2$, the distribution field rank $4$.

This is a perfectly good notation for writing down \emph{one} concrete shape.
The trouble starts the moment we try to write a \emph{function} that should
work at more than one rank. Consider the humble \code{axpy}
operation---``$\alpha$ times $\vx$ plus $\vy$''---the workhorse of every
iterative solver in this book:
\[
  \mathrm{axpy}(\alpha,\vx,\vy) \;=\; \alpha\,\vx + \vy .
\]
What does \code{axpy} require of its arguments? Only that $\vx$ and $\vy$ have
\emph{the same shape}, whatever that shape is, so that the scaled $\vx$ can be
added to $\vy$ entry by entry. It does not care whether they are vectors,
matrices, or rank-$7$ tensors. Scaling and adding are \emph{elementwise}: they
look at one entry at a time and never consult the rank at all.

Now watch what happens when we reach for our only shape notation so far and
try to write \code{axpy}'s signature. The natural-looking attempt is
\begin{lstlisting}[style=l4style]
axpy :: Scalar -> Tensor[N] -> Tensor[N] -> Tensor[N]   -- WRONG (pins rank 1)
\end{lstlisting}
The intent was surely ``the two tensors share a shape, call it \texttt{N}, and
the result has that shape too.'' But that is \emph{not} what this signature
says. The bracket \lstinline[style=l4style]|[N]| has exactly one entry, so it
denotes a tensor of \emph{rank one}: a flat vector of length \texttt{N}. By
writing \lstinline[style=l4style]|Tensor[N]| three times we have indeed said
``all three have the same shape''---but we have \emph{also} said, without
meaning to, ``and that shape is one-dimensional.'' A caller who hands
\code{axpy} two $3\times 3$ matrices will be turned away at the door, even
though scaling and adding $3\times 3$ matrices is exactly as sensible as
scaling and adding vectors.

\begin{pitfall}
\lstinline[style=l4style]|Tensor[N]| does \emph{not} mean ``the same shape as
the other operand.'' It means ``a tensor of rank one, with its single axis of
length \texttt{N}.'' The single letter inside the brackets is a name for
\emph{one axis}, not for \emph{a whole shape of unknown rank}. Reusing the
same letter across operands forces them to agree---but it forces them to agree
\emph{at rank one}. Using \lstinline[style=l4style]|Tensor[N]| to mean
``congruent, any rank'' is the single most common shape-notation error, and it
silently narrows a general operation to vectors only.
\end{pitfall}

The mismatch is worth seeing at a glance.
\Cref{fig:sh-pinned-vs-flex} sets the rank-$1$ commitment of
\lstinline[style=l4style]|Tensor[N]| beside the flexible notation we are about
to build. The cure is a way to name a \emph{whole run of axes}---possibly
several, possibly we-don't-know-how-many---and to say ``this other tensor has
the same run.'' That is the named shape group, and the rest of the chapter
develops it.

\begin{figure}[t]
\centering
\begin{tikzpicture}[scale=1.0]
  % --- LEFT: Tensor[N], one pinned box ---
  \begin{scope}[shift={(-3.6,0)}]
    \node[font=\scriptsize\bfseries, simRust] at (0,1.55) {\code{Tensor[N]}};
    \node[font=\scriptsize\itshape, simGray, text width=30mm, align=center]
      at (0,1.15) {one axis: rank pinned to 1};
    \draw[draw=simRust, fill=simBgRust, thick] (-0.5,-0.45) rectangle (0.5,0.45);
    \node[font=\small, simInk] at (0,0) {$N$};
    % a lock to convey "pinned"
    \node[font=\scriptsize, simRust] at (0,-0.95) {rank $=1$, fixed};
  \end{scope}
  % divider
  \draw[simGray!40, densely dotted] (-1.4,-1.4) -- (-1.4,1.7);
  % --- RIGHT: Tensor[(S: ...)], flexible bracket ---
  \begin{scope}[shift={(1.9,0)}]
    \node[font=\scriptsize\bfseries, simBlue] at (0.55,1.55)
      {\lstinline[style=l4style]|Tensor[(S: ...)]|};
    \node[font=\scriptsize\itshape, simGray, text width=42mm, align=center]
      at (0.55,1.15) {a named run \texttt{S}: any number of axes};
    % a row of boxes under a spanning bracket
    \foreach \i/\lab in {0/{},1/{},2/{}}{
      \draw[draw=simBlue, fill=simBgBlue, thick]
        (\i*0.95-0.45,-0.45) rectangle (\i*0.95+0.45,0.45);
      \node[font=\footnotesize, simGray] at (\i*0.95, 0) {$\cdot$};
    }
    \node[font=\scriptsize, simGray] at (3.05,0) {$\cdots$};
    \draw[decorate, decoration={brace, amplitude=5pt, mirror}, simBlue]
      (-0.55,-0.6) -- (2.45,-0.6)
      node[midway, below=6pt, font=\scriptsize, simBlue] {the group \texttt{S}};
  \end{scope}
\end{tikzpicture}
\caption[Pinned axis versus flexible group]{The rank-$1$ commitment of
\lstinline[style=l4style]|Tensor[N]| (left, rust) versus the rank-agnostic
named group \lstinline[style=l4style]|Tensor[(S: ...)]| (right, blue).
\lstinline[style=l4style]|Tensor[N]| is a single boxed axis: its rank is
locked at one. \lstinline[style=l4style]|Tensor[(S: ...)]| is a labelled
bracket that spans \emph{some} run of axes---one, three, however many---and
gives that run a name, \texttt{S}, without ever committing to how many there
are.}
\label{fig:sh-pinned-vs-flex}
\end{figure}

\section{Named shape groups}
\label{sec:sh-groups}

The fix is to give a \emph{name} to a run of axes without fixing how many
axes the run contains. We write the binding form
\[
  \texttt{(S: ...)}
\]
and read it ``the shape pattern \texttt{...}, called \texttt{S}.'' The three
dots \texttt{...} are a genuine piece of the notation---a \emph{rank
wildcard} that matches \emph{zero or more} axes of any extent. So
\lstinline[style=l4style]|Tensor[(S: ...)]| is ``a tensor whose entire shape
is some run of axes, and we will call that run \texttt{S}.'' It is true of a
vector (the run is one axis), of a matrix (two axes), of a rank-$7$ tensor
(seven axes), and even of a scalar (the run is empty). The notation has
declined to commit to a rank, which is exactly the freedom \code{axpy} needs.

A name is only useful if we can refer back to it. The second half of the
notation is the \emph{use} form, written with a dollar sigil:
\[
  \texttt{\$S}
\]
which means ``a run of axes \emph{congruent to} the group \texttt{S} bound
earlier''---the same number of axes, each of the same extent. The sigil is
the whole point: it makes a \emph{back-reference} look different from a
\emph{fresh binding}. \lstinline[style=l4style]|(S: ...)| introduces the name;
\lstinline[style=l4style]|$S| reuses it.

\begin{definition}[Named shape group]
\label{def:sh-group}
A \emph{named shape group} \lstinline[style=l4style]|(S: ...)| binds a name
\texttt{S} to a contiguous run of tensor axes \emph{without committing to the
rank or to the individual extents of that run}. The binding occurs exactly
once, at the group's first appearance. Every later occurrence
\lstinline[style=l4style]|$S|, written with the \texttt{\$} sigil, is a
\emph{use}: it asserts that the shape it labels is \emph{congruent} to the
bound run---the same number of axes, matched extent for extent. Two shapes
are \emph{congruent} when they are identical axis-for-axis (modulo the
shape algebra of \cref{sec:sh-equational}).
\end{definition}

With this in hand, \code{axpy}'s signature becomes what we always meant:
\begin{lstlisting}[style=l4style]
axpy :: Scalar -> Tensor[(S: ...)] -> Tensor[$S] -> Tensor[$S]
\end{lstlisting}
Read left to right: \code{axpy} takes a scalar; then a tensor of \emph{some}
shape, which binds the group \texttt{S}; then a second tensor whose shape must
be \emph{congruent} to \texttt{S}; and it returns a tensor, again congruent
to \texttt{S}. The first operand \emph{introduces} the shape; the second
operand and the result \emph{borrow} it. Nothing here mentions a rank, so the
one signature describes \code{axpy} on vectors, on matrices, and on every
higher tensor at once. \Cref{fig:sh-hero} is the picture to keep in mind: a
tensor is a row of axis-boxes, a group is a bracket spanning some of them, and
a \lstinline[style=l4style]|$S| use is an arrow pointing back to the binding,
asserting ``the same run, here, again.''

\begin{figure}[t]
\centering
\begin{tikzpicture}[scale=1.0]
  % Top: binding operand --- a tensor as a row of axis boxes, bracketed as (S: ...)
  \begin{scope}[shift={(0,1.5)}]
    \foreach \i in {0,1,2}{
      \draw[draw=simBlue, fill=simBgBlue, thick]
        (\i*0.9,-0.35) rectangle (\i*0.9+0.7,0.35);
    }
    \node[font=\scriptsize, simGray] at (2.55,0) {$\cdots$};
    \draw[draw=simBlue, fill=simBgBlue, thick] (3.0,-0.35) rectangle (3.7,0.35);
    \draw[decorate, decoration={brace, amplitude=4pt}, simBlue]
      (-0.05,0.5) -- (3.75,0.5)
      node[midway, above=4pt, font=\scriptsize\bfseries, simBlue] (Sbind)
      {binds \texttt{S}};
    \node[font=\scriptsize\ttfamily, simInk, left=4mm of {(-0.05,0)}]
      (xlab) at (-1.7,0) {\lstinline[style=l4style]|Tensor[(S: ...)]|};
  \end{scope}
  % Middle: use operand --- congruent, $S
  \begin{scope}[shift={(0,0)}]
    \foreach \i in {0,1,2}{
      \draw[draw=simTeal, fill=simBgGreen, thick]
        (\i*0.9,-0.35) rectangle (\i*0.9+0.7,0.35);
    }
    \node[font=\scriptsize, simGray] at (2.55,0) {$\cdots$};
    \draw[draw=simTeal, fill=simBgGreen, thick] (3.0,-0.35) rectangle (3.7,0.35);
    \draw[decorate, decoration={brace, amplitude=4pt, mirror}, simTeal]
      (-0.05,-0.5) -- (3.75,-0.5)
      node[midway, below=4pt, font=\scriptsize\bfseries, simTeal] (Suse)
      {uses \texttt{\$S}};
    \node[font=\scriptsize\ttfamily, simInk] at (-1.7,0)
      {\lstinline[style=l4style]|Tensor[$S]|};
  \end{scope}
  % the back-arrow: $S use --> S binding (congruence assertion)
  \draw[-{Stealth[length=2.6mm]}, simRust, thick]
    (4.0,-0.1) .. controls (4.9,0.55) and (4.9,1.45) .. (4.0,2.0)
    node[midway, right=1pt, font=\scriptsize\itshape, simRust, align=left]
    {congruent:\\same run};
\end{tikzpicture}
\caption[The shape-algebra hero figure]{A named shape group at work. Each
tensor is a row of axis-boxes; its rank is how many boxes there are, left
unspecified by the trailing \texttt{...}. The top operand's whole shape is
bracketed and \emph{bound} to the name \texttt{S}. The bottom operand is
written \lstinline[style=l4style]|Tensor[$S]|: the \texttt{\$} sigil makes it a
\emph{use}, and the rust back-arrow is the assertion it carries---``my shape is
\emph{congruent} to the run that bound \texttt{S}.'' Neither tensor's rank is
fixed; only their \emph{agreement} is.}
\label{fig:sh-hero}
\end{figure}

\subsection{Why not just reuse a shape variable?}
\label{sec:sh-vs-variable}

A sharp reader will object: we could have written \code{axpy} with a single
\emph{shape variable} $\sigma$ standing for the whole shape,
\begin{lstlisting}[style=l4style]
axpy :: Scalar -> Tensor[s] -> Tensor[s] -> Tensor[s]
\end{lstlisting}
where \texttt{s} (the typeset $\sigma$) names ``the entire shape, any rank.''
That is correct, and when the \emph{whole} shape is shared it is the lighter
spelling. Named groups earn their keep in the cases a bare variable cannot
reach, of which there are two.

First, a group can name a \emph{partial} run while leaving the rest free.
Suppose an operation shares a leading block of axes between two operands but
lets each carry its own trailing axis:
\begin{lstlisting}[style=l4style]
batched_op :: Tensor[(S: ...), k] -> Tensor[$S, m] -> Tensor[$S, k, m]
\end{lstlisting}
Here \texttt{S} is the shared leading run; \texttt{k} and \texttt{m} are
free trailing axes that differ between the operands. A single variable
$\sigma$ cannot express ``share \emph{these} axes but not \emph{that} one.''
The group splits a shape into a named congruent part and a free part---the
binding \lstinline[style=l4style]|(S: a, ...)| even mixes a pinned leading
axis \texttt{a} with a wildcard tail, so a group is as fine-grained as the
contract needs.

Second, and more quietly, a group makes the rank-agnosticism
\emph{visible}. A bare $\sigma$ relies on the reader knowing the convention
that a lone shape variable spans any rank; the trailing \texttt{...} inside
\lstinline[style=l4style]|(S: ...)| says it on the page. When the surrounding
code also uses concrete rank-$1$ shapes like
\lstinline[style=l4style]|Tensor[N]| for genuine flat vectors---and the next
section shows it does---the explicit wildcard is what stops a reader from
misreading the congruence contract as another rank-$1$ commitment.

\begin{keyidea}
A named shape group \lstinline[style=l4style]|(S: ...)| binds a name to a run
of axes \emph{of unknown rank}; a later use \lstinline[style=l4style]|$S|
asserts a second shape is \emph{congruent} to it. This is
\emph{rank-agnostic congruence}: one signature, e.g.\
\lstinline[style=l4style]|axpy :: Scalar -> Tensor[(S: ...)] -> Tensor[$S] -> Tensor[$S]|,
correctly describing an operation on vectors, on matrices, and on every
higher tensor simultaneously---because it constrains the operands to
\emph{agree} without ever constraining their \emph{rank}. The
\texttt{\$} sigil is the difference between a fresh binding and a back-reference.
\end{keyidea}

\section{Worked example: \texorpdfstring{\code{axpy}}{axpy} at two ranks}
\label{sec:sh-wex-axpy}

The claim is that one signature serves all ranks. The way to believe it is to
run the same signature against two concrete arguments---a vector and a
matrix---and watch the congruence check succeed both times, then watch the
wrong signature reject the matrix.

\begin{workedexample}{Why \code{axpy} wants a named group}{sh-axpy}
\textbf{The signature under test.}
\begin{lstlisting}[style=l4style]
axpy :: Scalar -> Tensor[(S: ...)] -> Tensor[$S] -> Tensor[$S]
\end{lstlisting}
We instantiate it twice and resolve the group \texttt{S} each time.

\tcblower
\textbf{Instance 1 --- a rank-$1$ vector.} Apply \code{axpy} with
$\alpha = 2$, $\vx,\vy$ of shape \lstinline[style=l4style]|[3]| (length-$3$
vectors). The first tensor operand binds the group:
\[
  \texttt{S} \;\mapsto\; \texttt{[3]} \qquad(\text{rank } 1).
\]
The second operand's required shape is \lstinline[style=l4style]|$S = [3]|: the
supplied $\vy$ is \lstinline[style=l4style]|[3]|, which is congruent.
\emph{Check passes.} The result shape is \lstinline[style=l4style]|$S = [3]|,
a length-$3$ vector. Elementwise: $2x_i + y_i$ at each of the $3$ entries.

\textbf{Instance 2 --- a rank-$2$ matrix-shaped field.} Now apply the
\emph{same} \code{axpy} with $\alpha = 2$, but $\vx,\vy$ of shape
\lstinline[style=l4style]|[3, 3]| (think of a $3\times 3$ block of a
tensor-valued material coefficient, or a matrix-shaped field). The first
operand binds the group to a \emph{rank-$2$} run:
\[
  \texttt{S} \;\mapsto\; \texttt{[3, 3]} \qquad(\text{rank } 2).
\]
The second operand's required shape is now \lstinline[style=l4style]|$S = [3, 3]|;
the supplied $\vy$ is \lstinline[style=l4style]|[3, 3]|, congruent.
\emph{Check passes.} The result is \lstinline[style=l4style]|[3, 3]|;
elementwise $2x_{ij} + y_{ij}$ across all $9$ entries. The signature never
mentioned a rank, so it bound \texttt{S} to whatever the first operand
brought---rank $1$ in instance~1, rank $2$ in instance~2.

\textbf{The wrong signature, same two instances.} Replace the group with the
rank-$1$ form:
\begin{lstlisting}[style=l4style]
axpy_bad :: Scalar -> Tensor[N] -> Tensor[N] -> Tensor[N]
\end{lstlisting}
\emph{Instance 1} still type-checks: the vectors are
\lstinline[style=l4style]|[3]|, so \texttt{N} resolves to $3$ and all three
slots agree. \emph{Instance 2 fails.} The matrices are
\lstinline[style=l4style]|[3, 3]|---rank $2$---but
\lstinline[style=l4style]|Tensor[N]| demands rank $1$. There is no value of
the single axis \texttt{N} that can match a two-axis shape: matching
\lstinline[style=l4style]|[N]| against \lstinline[style=l4style]|[3, 3]|
fails on \emph{rank} before any extent is even compared. The matrix case is
rejected---not because scaling-and-adding a matrix is ill-defined (it is
perfectly well-defined) but because the signature accidentally outlawed it.
\end{workedexample}

\begin{remark}
Notice \emph{where} the two signatures diverge: at the rank check, which
happens \emph{before} extents are compared. \lstinline[style=l4style]|Tensor[N]|
fixes the number of axes to one; \lstinline[style=l4style]|Tensor[(S: ...)]|
leaves it open and binds it from the argument. Everything downstream---the
elementwise arithmetic, the laws like
\lstinline[style=l4style]|axpy 0 v y = y|---is identical between the two ranks.
The \emph{only} thing the group buys is permission to run at a rank the
operation never cared about in the first place.
\end{remark}

\Cref{fig:sh-boundary} abstracts what just happened into the rule the matcher
follows at \emph{every} function boundary: bind the group from the first
shape it labels, then test each later use for congruence, rank first.

\begin{figure}[t]
\centering
\begin{tikzpicture}[node distance=7mm and 10mm]
  % the function boundary as a vertical gate
  \node[stage, minimum height=30mm, text width=4mm] (gate) at (0,0) {};
  \node[font=\scriptsize\bfseries, simBlue, above=1mm of gate]
    {function boundary};
  % incoming binding operand
  \node[data, left=14mm of gate.west, yshift=8mm, text width=22mm] (bind)
    {1st operand\\\lstinline[style=l4style]|Tensor[(S: ...)]|};
  % incoming use operand
  \node[data, left=14mm of gate.west, yshift=-8mm, text width=22mm] (use)
    {2nd operand\\\lstinline[style=l4style]|Tensor[$S]|};
  % the resolved value of S (the gate's memory)
  \node[font=\scriptsize\ttfamily, simRust, right=4mm of gate, yshift=8mm,
        align=left] (resolved)
    {bind:\\ $\texttt{S}\mapsto$ this shape};
  \node[font=\scriptsize, simTeal, right=4mm of gate, yshift=-8mm,
        align=left] (checked)
    {check:\\ rank? then extents?};
  \draw[flow] (bind.east) -- node[above, font=\scriptsize, simGray]{bind} (gate.west|-bind.east);
  \draw[flow] (use.east) -- node[below, font=\scriptsize, simGray]{test} (gate.west|-use.east);
  \draw[refedge, simRust] (gate.east|-resolved) -- (resolved);
  \draw[refedge, simTeal] (gate.east|-checked) -- (checked);
  % the two-step rule annotated below
  \node[font=\scriptsize\itshape, simGray, below=8mm of gate, text width=70mm,
        align=center]
    {\textbf{1.}~rank mismatch $\Rightarrow$ reject \emph{immediately}
     (this is where \lstinline[style=l4style]|Tensor[N]| fails the matrix).\quad
     \textbf{2.}~ranks agree $\Rightarrow$ compare extents, axis by axis,
     modulo the shape algebra (\cref{sec:sh-equational}).};
\end{tikzpicture}
\caption[The congruence check at a boundary]{What the shape matcher does at a
function boundary. The \emph{first} shape that labels a group \texttt{S}
\emph{binds} it (rust)---fixing both the rank and the extents \texttt{S} now
stands for. Every \emph{later} use \lstinline[style=l4style]|$S| is
\emph{tested} for congruence (teal): \textbf{first} the rank must match---a
mismatch is rejected on the spot, which is exactly how
\lstinline[style=l4style]|Tensor[N]| turns away a rank-$2$ argument
(\cref{wex:sh-axpy})---and \textbf{then}, only if the ranks agree, the
extents are compared axis by axis. The same two-step gate guards \code{axpy},
\code{apply\_linop}, and every other signature in this book.}
\label{fig:sh-boundary}
\end{figure}

\section{Operator shapes: domain and range}
\label{sec:sh-operators}

Congruence groups describe operations whose operands share a shape. But the
most important objects in this book---the linear operators $\mat{K}$, $\mat{M}$,
$\mat{A}(\omega)$ of \cref{ch:targets}---are maps whose \emph{input} and
\emph{output} shapes generally \emph{differ}. A matrix that takes a length-$n$
vector to a length-$m$ vector has two shapes to name, not one. The notation
extends to cover this with two groups: a \emph{range} group and a
\emph{domain} group.

We write them \emph{range-first}:
\begin{lstlisting}[style=l4style]
LinOp[(R: ...), (D: ...)]
\end{lstlisting}
for an operator from domain shape \texttt{D} to range shape \texttt{R}, both
of unknown rank. The order---range, then domain---is the matrix convention in
disguise. An $m\times n$ matrix is written rows-first ($m$) and maps a
length-$n$ \emph{domain} to a length-$m$ \emph{range}; the symbol order
$(m,n)$ is (range, domain). \lstinline[style=l4style]|LinOp[(R: ...), (D: ...)]|
is the rank-agnostic version of that same $(m,n)$ ordering. Applying the
operator is then a pair of uses, one for each group:
\begin{lstlisting}[style=l4style]
apply_linop :: LinOp[(R: ...), (D: ...)] -> Tensor[$D] -> Tensor[$R]
\end{lstlisting}
the operand is a use of the \emph{domain} group (\lstinline[style=l4style]|$D|);
the result is a use of the \emph{range} group (\lstinline[style=l4style]|$R|).
\Cref{fig:sh-operator} draws the flow: a domain-shaped tensor enters, the
operator carries it across, a range-shaped tensor leaves.

\begin{figure}[t]
\centering
\begin{tikzpicture}[node distance=10mm and 13mm]
  \node[data, text width=18mm] (in) {input\\\lstinline[style=l4style]|Tensor[$D]|};
  \node[stage, right=of in, text width=34mm] (op)
    {\lstinline[style=l4style]|LinOp[(R: ...), (D: ...)]|};
  \node[data, right=of op, text width=18mm] (out)
    {output\\\lstinline[style=l4style]|Tensor[$R]|};
  \draw[flow] (in) -- node[above, font=\scriptsize, simGray] {domain} (op);
  \draw[flow] (op) -- node[above, font=\scriptsize, simGray] {range} (out);
  % range-first annotation under the operator
  \node[font=\scriptsize\itshape, simGray, below=5mm of op, text width=46mm,
        align=center]
    {written \emph{range-first}: \texttt{(R: ...)} then \texttt{(D: ...)},\\
     matching the $m\times n$ (rows $\times$ cols) matrix order};
  % brace pointing out the (R,D) ordering
  \draw[decorate, decoration={brace, amplitude=4pt}, simBlue]
    ([yshift=6mm]op.north west) -- ([yshift=6mm]op.north east)
    node[midway, above=4pt, font=\scriptsize, simBlue]
    {range $\leftarrow$ domain};
\end{tikzpicture}
\caption[Operator shapes are range-first]{A linear operator names two shape
groups: a range group \texttt{R} and a domain group \texttt{D}, written
\emph{range-first} to match the rows-then-columns ordering of an $m\times n$
matrix. Application \lstinline[style=l4style]|apply_linop op| takes a
domain-shaped input \lstinline[style=l4style]|Tensor[$D]| and produces a
range-shaped output \lstinline[style=l4style]|Tensor[$R]|. Both groups are
rank-agnostic: the operator could map vectors to vectors, fields to fields,
or one rank to another.}
\label{fig:sh-operator}
\end{figure}

\subsection{The square case: \texorpdfstring{$\mat{A}\colon V\to V$}{A : V to V}}
\label{sec:sh-square}

A great many operators in this book are \emph{endomorphisms}: their domain and
range are the same space. The system matrix $\mat{K}$ of an eigenproblem, a
preconditioner $\mat{M}^{-1}$, the time-stepping operator of a transient
solve---each takes a state vector and returns a vector of the very same shape.
For these we do not need two independent groups; we bind \emph{one} and use it
for the range as well:
\begin{lstlisting}[style=l4style]
LinOp[(S: ...), $S]
\end{lstlisting}
``an operator whose range shape \emph{is} its domain shape \texttt{S}.'' This
is the operator-shape echo of the congruent \code{axpy} signature: there, the
two \emph{operands} shared a group; here, the operator's \emph{domain and
range} share one. Applied, a square operator threads a shape straight through:
\begin{lstlisting}[style=l4style]
apply_linop :: LinOp[(S: ...), $S] -> Tensor[$S] -> Tensor[$S]
\end{lstlisting}
which is exactly the signature a Krylov solver wants from its operator
(\cref{ch:matrixfree})---feed it a vector, get back a vector you can feed it
again.

\begin{workedexample}{An operator-apply congruence check}{sh-matvec}
\textbf{Setup.} Take the matrix-free mass operator $\mat{M}$ of
\cref{sec:mf-wex1}, acting on the flat global dof-vector. At L1 the global
vector is a genuine rank-$1$ object, so we instantiate the square form with a
concrete rank-$1$ domain.
\begin{lstlisting}[style=l4style]
applyM :: LinOp[(S: ...), $S] -> Tensor[$S] -> Tensor[$S]
\end{lstlisting}

\tcblower
\textbf{Bind from the operator's domain.} The operator $\mat{M}$ acts on
length-$N$ dof-vectors, $N=3$ in that example. Instantiating the square
operator type against $\mat{M}$ resolves
\[
  \texttt{S} \;\mapsto\; \texttt{[3]} ,
\]
so both the domain and range are \lstinline[style=l4style]|[3]|.

\textbf{Check the operand.} We apply $\mat{M}$ to $\vv=(1,2,3)$, shape
\lstinline[style=l4style]|[3]|. The required operand shape is
\lstinline[style=l4style]|$S = [3]|; the supplied $\vv$ is
\lstinline[style=l4style]|[3]|. \emph{Congruent---check passes.}

\textbf{Check the result.} The result shape is \lstinline[style=l4style]|$S = [3]|.
The pipeline of \cref{sec:mf-wex1} returned $\mat{M}\vv = (0.667, 2.000, 1.333)$,
shape \lstinline[style=l4style]|[3]|. \emph{Congruent---and crucially, the
output can be fed straight back into \code{applyM}}, because its shape is again
\texttt{S}. That re-feedability is precisely what the square form
\lstinline[style=l4style]|LinOp[(S: ...), $S]| guarantees, and it is what makes
the Krylov subspace $\{\vv, \mat{M}\vv, \mat{M}^2\vv, \dots\}$ of
\cref{ch:matrixfree} type-check at every power.
\end{workedexample}

\begin{notationbox}
At the \emph{lowest} layers of the spec---where Palace's operators act on
genuine flat dof-vectors of a definite length---the older concrete spelling
\lstinline[style=l4style]|LinearOperator[M, N]| over rank-$1$
\lstinline[style=l4style]|Tensor[N]| vectors is faithful and is kept. There
the lengths $M$, $N$ really are single flat-vector lengths, and the rank-$1$
commitment is correct rather than accidental. The group form
\lstinline[style=l4style]|LinOp[(R: ...), (D: ...)]| is the rank-agnostic
\emph{calculus} rendering, used at the higher layers where the same operator
is described over structured, possibly higher-rank shapes. Both spellings
coexist; choosing between them is choosing whether the rank is genuinely
fixed (keep \lstinline[style=l4style]|[N]|) or genuinely open (use a group).
\end{notationbox}

\section{Named axes of fixed meaning}
\label{sec:sh-fixedaxes}

There is a second kind of name in our shape brackets, and it is easy to
confuse with a congruence group until the distinction is drawn sharply. A
congruence group like \texttt{S} is a \emph{variable}: it stands for whatever
run of axes the first operand happens to bring, and its whole job is to assert
that two shapes \emph{agree}. A \emph{named axis of fixed meaning} is the
opposite: it is a single axis whose letter denotes one specific, fixed
quantity, the same way \lstinline[style=l4style]|Tensor[H, W, VY=3, VX=3]| of
\cref{ch:bigidea} names a height, a width, and two velocity axes. \texttt{H}
is not a congruence variable to be resolved; it is \emph{the height}, always.

The book's recurring family of fixed-meaning axes is the element-local
vocabulary of the matrix-free kernel, which we met in \cref{ch:matrixfree}.
Recall its five letters:
\begin{description}[leftmargin=1.6em, style=nextline, nosep]
\item[\texttt{E}] the element count---how many mesh elements the operator runs over.
\item[\texttt{L}] local dofs per element, $(p+1)^d$ for a degree-$p$ element.
\item[\texttt{P}] quadrature points per element.
\item[\texttt{C}] value (or derivative) components carried at a quadrature point.
\item[\texttt{G}] geometry-factor components stored per quadrature point.
\end{description}
and its three rank-structured tensors, \lstinline[style=l4style]|Tensor[(E, L)]|,
\lstinline[style=l4style]|Tensor[(E, P, C)]|, and
\lstinline[style=l4style]|Tensor[(E, P, G)]|. Each of these five letters has a
\emph{definite} meaning fixed by the finite-element construction; none of them
is a back-reference, and none is rank-agnostic.
\Cref{fig:sh-axis-vs-group} sets the two kinds of name side by side.

\begin{figure}[t]
\centering
\begin{tikzpicture}[scale=1.0]
  % LEFT: fixed-meaning axes E, P, C
  \begin{scope}[shift={(-3.5,0)}]
    \node[font=\scriptsize\bfseries, simGreen] at (0.0,1.7)
      {\lstinline[style=l4style]|Tensor[(E, P, C)]|};
    \node[font=\scriptsize\itshape, simGray, text width=34mm, align=center]
      at (0.0,1.25) {fixed-meaning axes};
    \foreach \i/\lab/\mean in {0/E/{elements},1/P/{quad pts},2/C/{comps}}{
      \draw[draw=simGreen, fill=simBgGreen, thick]
        (\i*1.05-0.45,-0.35) rectangle (\i*1.05+0.45,0.35);
      \node[font=\small, simInk] at (\i*1.05,0) {$\lab$};
      \node[font=\scriptsize, simGray] at (\i*1.05,-0.75) {\mean};
    }
    \node[font=\scriptsize\itshape, simGreen, text width=34mm, align=center]
      at (0.0,-1.35) {each letter = one fixed quantity\\ (a definite role)};
  \end{scope}
  % divider
  \draw[simGray!40, densely dotted] (-1.05,-1.7) -- (-1.05,2.0);
  % RIGHT: congruence group S
  \begin{scope}[shift={(2.0,0)}]
    \node[font=\scriptsize\bfseries, simBlue] at (0.55,1.7)
      {\lstinline[style=l4style]|Tensor[(S: ...)] / Tensor[$S]|};
    \node[font=\scriptsize\itshape, simGray, text width=42mm, align=center]
      at (0.55,1.25) {congruence variable};
    \foreach \i in {0,1}{
      \draw[draw=simBlue, fill=simBgBlue, thick]
        (\i*0.85,-0.35) rectangle (\i*0.85+0.65,0.35);
    }
    \node[font=\scriptsize, simGray] at (1.95,0) {$\cdots$};
    \draw[decorate, decoration={brace, amplitude=4pt, mirror}, simBlue]
      (-0.05,-0.5) -- (2.3,-0.5)
      node[midway, below=4pt, font=\scriptsize, simBlue] {\texttt{S}: any run};
    \node[font=\scriptsize\itshape, simBlue, text width=42mm, align=center]
      at (0.55,-1.35) {one name = ``the same run again''\\ (a back-reference, any rank)};
  \end{scope}
\end{tikzpicture}
\caption[Fixed-meaning axis versus congruence group]{Two kinds of name in a
shape bracket. \emph{Left (green):} a \emph{fixed-meaning axis}---\texttt{E},
\texttt{P}, \texttt{C} of the matrix-free kernel (\cref{ch:matrixfree}). Each
letter denotes one specific finite-element quantity; the rank is fixed (here,
three) and each axis has a definite role. \emph{Right (blue):} a
\emph{congruence group} \texttt{S}. The single name does not denote a quantity
at all; it denotes ``the same run of axes as wherever \texttt{S} was bound,''
of any rank. A fixed-meaning axis answers \emph{what is this axis};
a congruence group answers \emph{does this shape match that one}.}
\label{fig:sh-axis-vs-group}
\end{figure}

\begin{keyidea}
Distinguish two kinds of name in a shape bracket. A \emph{named axis of fixed
meaning} (\texttt{E}, \texttt{L}, \texttt{P}, \texttt{C}, \texttt{G};
\texttt{H}, \texttt{W}) is \emph{one} axis with a definite, permanent role---it
answers ``what is this axis.'' A \emph{named shape group}
(\texttt{(S: ...)} / \texttt{\$S}) is a \emph{congruence variable} spanning a
run of \emph{unknown} rank---it answers ``does this shape match that one.''
The first is concrete and rank-bearing; the second is abstract and
rank-agnostic. They share the look of a letter in brackets and nothing else.
\end{keyidea}

The two kinds even meet at a boundary. The flat global dof-vector is a genuine
rank-$1$ \lstinline[style=l4style]|Tensor[N]|---neither a fixed-meaning
element-local axis nor a congruence group, just a flat vector. The gather and
scatter of \cref{ch:matrixfree} are exactly the crossing between that flat
\texttt{N} world and the fixed-meaning \lstinline[style=l4style]|(E, L)| world:
\lstinline[style=l4style]|element_restrict :: ... -> Tensor[N] -> Tensor[(E, L)]|
turns one flat vector into the element-local family, and its transpose turns it
back. Three different notational species---a rank-$1$ flat axis, a congruence
group, a fixed-meaning family---live in the same pipeline, and keeping them
straight is what makes the pipeline legible.

\section{Scalar promotion: one operator, two scalar types}
\label{sec:sh-promotion}

So far we have reasoned about \emph{shapes}. There is one rule about
\emph{element types} that the shape contracts lean on, and it is worth stating
in its own right because it removes what would otherwise be a swarm of
near-duplicate operators. It is called \emph{scalar promotion}.

The setting: Palace works with both \emph{real} and \emph{complex} tensors.
The driven and eigenmode analyses (\cref{ch:targets}) produce complex fields;
the electrostatic and magnetostatic analyses are real. An operation like
\code{axpy} therefore comes in a real flavour (real scalar, real vectors) and
a complex flavour (complex scalar, complex vectors). But there is a third
combination that arises constantly in practice: a \emph{real} scalar
multiplying a \emph{complex} vector. A step size, a relaxation weight, a
quadrature factor---these are real numbers, and they routinely scale complex
fields.

Rather than mint a separate operator for that mixed case, the calculus admits
a single typing rule:
\[
  \boxed{\ \texttt{real} \;\sqsubseteq\; \texttt{complex}\ }
\]
read ``a real scalar may stand wherever a complex scalar is expected,''
promoted to a complex number with \emph{zero imaginary part}. The promotion is
\emph{exact}: zero is representable exactly, so $\alpha$ and $\alpha + 0i$ are
the very same number to the last bit, and every algebraic law of the operator
holds identically before and after promotion. \Cref{fig:sh-lattice} draws the
one-step lattice.

\begin{figure}[t]
\centering
\begin{tikzpicture}[scale=1.0]
  \node[draw=simViolet, fill=simBgViolet, rounded corners=3pt, thick,
        minimum width=20mm, minimum height=8mm, font=\small] (c) at (0,1.4)
    {\texttt{complex}};
  \node[draw=simBlue, fill=simBgBlue, rounded corners=3pt, thick,
        minimum width=20mm, minimum height=8mm, font=\small] (r) at (0,0)
    {\texttt{real}};
  \draw[-{Stealth[length=2.6mm]}, simRust, thick] (r) -- (c)
    node[midway, right=2mm, font=\scriptsize\itshape, simRust, align=left]
    {promote:\\ $\alpha \mapsto \alpha + 0i$\\ (exact)};
  \node[font=\scriptsize\itshape, simGray, left=12mm of {(0,0.7)}, text width=26mm,
        align=right] at (-2.1,0.7)
    {the lattice\\ \texttt{real} $\sqsubseteq$ \texttt{complex}};
  % the forbidden reverse, struck out
  \draw[-{Stealth[length=2.2mm]}, simGray, dotted] (c.west) to[bend right=40]
    node[midway, left=1pt, font=\scriptsize, simGray] {\,no} (r.west);
\end{tikzpicture}
\caption[The scalar-promotion lattice]{Scalar promotion is the one-step
ordering \texttt{real} $\sqsubseteq$ \texttt{complex}. A real scalar may be
\emph{promoted} upward into a complex operation by attaching a zero imaginary
part, $\alpha \mapsto \alpha + 0i$---exactly, with no rounding. The reverse
(grey, dotted) is forbidden: a complex scalar cannot be demoted into a real
operation, because its imaginary part would have nowhere to go. The promotion
collapses what would otherwise be two near-identical operators (real-scalar
and complex-scalar overloads) into one.}
\label{fig:sh-lattice}
\end{figure}

\begin{definition}[Scalar promotion]
\label{def:sh-promotion}
For a tensor operator with scalar arguments and tensor operands of element
type $T \in \{\texttt{real}, \texttt{complex}\}$:
if $T = \texttt{complex}$, a scalar argument supplied as \texttt{real} is
\emph{promoted} to \texttt{complex} by setting its imaginary part to zero, and
the operator is applied as its complex form. The promotion is exact and the
operator's algebraic laws are invariant under it. It is a property of the
\emph{type system}---a coercion at the boundary---not a separate operator.
\end{definition}

This is the difference between a \emph{typing rule} and an \emph{operator
variant}. Were the real-scalar-against-complex-vector case a distinct
operator, every algebraic law would have to be restated for it, every call
site would have to choose which overload it meant, and every lowering would
carry the real fast-path as algebraic content. Promoting instead collapses all
of that: there is one \code{axpy}, its scalar argument's type is lifted by
\texttt{real} $\sqsubseteq$ \texttt{complex} when the vectors are complex, and
the laws are written once.

\begin{workedexample}{A real coefficient times a complex field}{sh-promotion}
\textbf{Setup.} A driven-analysis update scales a complex field $\vx$ by a
real relaxation weight and adds it to a complex accumulator $\vy$:
\[
  \alpha = 0.5 \ (\texttt{real}), \qquad
  \vx = (1+2i,\; 3-i), \qquad
  \vy = (i,\; 2),
\]
both vectors \lstinline[style=l4style]|Tensor[2]| of \texttt{complex} element
type. We call \lstinline[style=l4style]|axpy alpha x y|.

\tcblower
\textbf{Shape side (the group).} The signature
\lstinline[style=l4style]|axpy :: Scalar -> Tensor[(S: ...)] -> Tensor[$S] -> Tensor[$S]|
binds $\texttt{S} \mapsto \texttt{[2]}$ from $\vx$; $\vy$ and the result are
congruent. \emph{Shape check passes}---this is the same machinery as
\cref{wex:sh-axpy}, indifferent to the element type.

\textbf{Type side (the promotion).} The vectors are \texttt{complex}, but
$\alpha = 0.5$ is \texttt{real}. By \texttt{real} $\sqsubseteq$ \texttt{complex}
it is promoted to $0.5 + 0i$, \emph{exactly}. No separate ``real-times-complex
axpy'' is invoked; the one complex \code{axpy} runs with a promoted scalar.

\textbf{Compute.}
\begin{align*}
  \alpha\vx + \vy
  &= (0.5{+}0i)\,(1{+}2i,\; 3{-}i) + (i,\; 2) \\
  &= (0.5{+}i,\; 1.5{-}0.5i) + (i,\; 2)
   = (0.5 + 2i,\; 3.5 - 0.5i).
\end{align*}
The result is \lstinline[style=l4style]|Tensor[2]| complex, congruent to
\texttt{S}, as the signature promised.
\end{workedexample}

\subsection{Where promotion does \emph{not} apply}
\label{sec:sh-promotion-limits}

The rule is deliberately narrow, and its edges matter as much as its content.

\begin{itemize}[leftmargin=1.4em]
\item \textbf{Not the reverse direction.} A \emph{complex} scalar against a
\emph{real} tensor is \emph{not} promoted---there is no demotion. The complex
scalar's imaginary part has no destination in a real result, so the lattice
runs one way only (\cref{fig:sh-lattice}, dotted arrow). This combination is a
type error, not a coercion.

\item \textbf{Not reductions whose output is already a scalar.} An operation
like \lstinline[style=l4style]|dot :: Tensor[$S] -> Tensor[$S] -> Scalar| has
no \emph{input} scalar to promote---its scalar is the \emph{output}, and the
output type is fixed by the vectors' element type (complex vectors give a
complex result). The same goes for a norm
\lstinline[style=l4style]|nrm2 :: Tensor[$S] -> Real|, whose result is real
regardless. Promotion is about an \emph{incoming} scalar argument; a reduction
has none.

\item \textbf{All-or-none across a scalar tuple.} An operator with several
scalar arguments (such as
\lstinline[style=l4style]|axpby :: Scalar -> Tensor[$S] -> Scalar -> Tensor[$S] -> Tensor[$S]|)
promotes \emph{all} of its scalars together or none. A mixture---one real, one
complex, against complex vectors---is not part of the contract.
\end{itemize}

\begin{pitfall}
Scalar promotion lifts a \texttt{real} scalar \emph{into} a complex operation
($\alpha \mapsto \alpha + 0i$); it never lifts a complex scalar into a real
one, and it never touches a reduction's \emph{output} scalar. If you find
yourself wanting to ``promote'' a complex coefficient down to a real result,
the type is genuinely wrong---there is no information-preserving way to do it,
and the calculus rejects it rather than silently discarding the imaginary part.
\end{pitfall}

\section{The shape equational theory, lightly}
\label{sec:sh-equational}

One loose end remains in \cref{def:sh-group}: congruence asks that two shapes
be ``identical axis-for-axis''---but identical in what sense? Sometimes two
shapes that look different on the page denote the same run of extents. A
reshape that merges a height axis \texttt{h} and a width axis \texttt{w} into a
single flattened axis produces an extent written \lstinline[style=l4style]|h*w|;
an operation that consumes the flattened tensor must recognise
\lstinline[style=l4style]|h*w| and \lstinline[style=l4style]|w*h| as the same
length, and \lstinline[style=l4style]|[h, w]| as reshape-compatible with
\lstinline[style=l4style]|[h*w]|. Shape matching is therefore not raw
character-by-character comparison; it is comparison \emph{modulo a small
algebra} of axis extents---the obvious arithmetic, commutative and associative
in $+$ and $\times$.

For the working purposes of this chapter, two facts suffice. First, when a
signature reuses a group (\lstinline[style=l4style]|$S| after
\lstinline[style=l4style]|(S: ...)|) or repeats a shape variable, the matcher
must (i) agree on \emph{rank}---the same number of axes---and then (ii) agree
on each axis extent \emph{up to this small algebra}, so that
\lstinline[style=l4style]|h*w| matches \lstinline[style=l4style]|w*h|. Second,
the rank check comes first: a rank mismatch fails immediately, which is exactly
why \lstinline[style=l4style]|Tensor[N]| could never match a rank-$2$ argument
in \cref{wex:sh-axpy}, regardless of any extent arithmetic. The full,
formal treatment of this shape algebra---the rewriting system that decides when
two extent expressions are equal---is deferred to \cref{app:shapes}; the
intuition above is all the running text needs.

\begin{notationbox}
A practical consequence: prefer to write a shared shape as a \emph{group} or a
\emph{variable}, not as an algebraic coincidence. If two operands must agree,
say so with \lstinline[style=l4style]|(S: ...)| / \lstinline[style=l4style]|$S|
and let the matcher enforce it. Reserve the extent algebra
(\lstinline[style=l4style]|h*w| and friends) for the places where a genuine
reshape or factorisation has \emph{produced} a composite extent. Congruence is
a contract you assert; the extent algebra is a fact the matcher verifies.
\end{notationbox}

\summarysection
A tensor's \emph{rank} is the number of axes in its shape bracket, and the
naive notation \lstinline[style=l4style]|Tensor[N]| pins that rank to one: it
names a single axis, so reusing it across operands forces them to agree
\emph{at rank one}, silently outlawing the matrix and higher-rank cases that
elementwise operations handle perfectly well. The cure is the \emph{named
shape group}: \lstinline[style=l4style]|(S: ...)| binds a name to a run of axes
of \emph{unknown} rank, and a later use \lstinline[style=l4style]|$S| (the
\texttt{\$} sigil marking a back-reference) asserts a second shape is
\emph{congruent} to it. This gives \emph{rank-agnostic congruence}: one
signature, \lstinline[style=l4style]|axpy :: Scalar -> Tensor[(S: ...)] -> Tensor[$S] -> Tensor[$S]|,
correct on vectors, matrices, and every higher tensor at once. Operators whose
domain and range shapes differ name two groups \emph{range-first},
\lstinline[style=l4style]|LinOp[(R: ...), (D: ...)]|, collapsing to the square
form \lstinline[style=l4style]|LinOp[(S: ...), $S]| for endomorphisms.
Congruence groups are \emph{variables} and must be kept distinct from
\emph{named axes of fixed meaning}---the \texttt{E}/\texttt{L}/\texttt{P}/%
\texttt{C}/\texttt{G} family of the matrix-free kernel---whose letters denote
definite, rank-bearing quantities. Finally, \emph{scalar promotion}
(\texttt{real} $\sqsubseteq$ \texttt{complex}) is a typing rule, not an
operator variant: a real scalar joins a complex operation by gaining a zero
imaginary part, exactly, collapsing the real/complex overloads into one
operator---but never in reverse, and never on a reduction's output scalar.

\furtherreading
The view of a matrix as a map between shaped spaces, rather than a grid of
numbers, is the organising idea of Trefethen and
Bau~\citep{trefethenbau1997}; their opening lectures on matrix--vector
products and the four fundamental subspaces are the linear-algebra backdrop for
the range-first operator-shape convention of \cref{sec:sh-operators}. The
named-axis style of shape annotation---carrying an axis's \emph{meaning} in its
name---is folklore in the array-programming and tensor-compiler communities;
the matrix-free element-local family of \cref{sec:sh-fixedaxes} is a concrete
instance, developed in full in \cref{ch:matrixfree}. The formal shape algebra
underlying \cref{sec:sh-equational} is collected in \cref{app:shapes}.

\exercisessection
\begin{enumerate}[nosep]
  \item \textbf{Rank counting.} For each shape, state the rank and say whether
  it could be the binding of a group \lstinline[style=l4style]|(S: ...)| in a
  call where the other operand is a length-$4$ vector:
  \lstinline[style=l4style]|[4]|, \lstinline[style=l4style]|[2, 2]|,
  \lstinline[style=l4style]|[4, 1]|, \lstinline[style=l4style]|[]|.
  (Recall congruence requires agreement on rank \emph{first}.)

  \item \textbf{The wrong signature, again.} A colleague writes a scaling
  operator as \lstinline[style=l4style]|scale :: Scalar -> Tensor[M] -> Tensor[M]|
  and reports it ``works fine on all our test vectors.'' Construct the smallest
  input on which it wrongly fails, and rewrite the signature so it does not.

  \item \textbf{Partial congruence.} Write a signature for an operation that
  takes two tensors sharing a leading run of axes but each carrying its own
  \emph{single} trailing axis (lengths \texttt{p} and \texttt{q} respectively),
  and returns a tensor with the shared run followed by \emph{both} trailing
  axes. (Hint: \lstinline[style=l4style]|(S: ...)| with two free trailing
  dims.)

  \item \textbf{Range-first.} An operator maps a rank-$2$ field of shape
  \lstinline[style=l4style]|[a, b]| to a rank-$1$ vector of shape
  \lstinline[style=l4style]|[c]|. Write its
  \lstinline[style=l4style]|LinOp[(R: ...), (D: ...)]| type and the type of its
  \lstinline[style=l4style]|apply_linop|. Which group is which?

  \item \textbf{Square or not.} Decide for each whether the operator is an
  endomorphism (deserving \lstinline[style=l4style]|LinOp[(S: ...), $S]|) or a
  general map (deserving two groups): (a) a preconditioner $\mat{M}^{-1}$;
  (b) the gather \lstinline[style=l4style]|element_restrict| of
  \cref{ch:matrixfree}; (c) the time-stepping operator of a transient solve;
  (d) the basis evaluation $B_{\diffop}$ mapping
  \lstinline[style=l4style]|[E, L]| to \lstinline[style=l4style]|[E, P, C]|.

  \item \textbf{Axis or group?} For each letter as it appears in the named
  shape, classify it as a \emph{fixed-meaning axis} or a \emph{congruence
  group}, and justify: \texttt{E} in \lstinline[style=l4style]|Tensor[(E, L)]|;
  \texttt{S} in \lstinline[style=l4style]|Tensor[$S]|; \texttt{H} in
  \lstinline[style=l4style]|Tensor[H, W, VY=3, VX=3]|; \texttt{D} in
  \lstinline[style=l4style]|LinOp[(R: ...), (D: ...)]|.

  \item \textbf{Promotion or error?} For each call against \texttt{complex}
  vectors $\vx,\vy$ unless stated, say whether scalar promotion applies, does
  nothing, or is a type error: (a) \lstinline[style=l4style]|axpy| with a real
  $\alpha$; (b) \lstinline[style=l4style]|axpy| with a complex $\alpha$ against
  \emph{real} $\vx,\vy$; (c) \lstinline[style=l4style]|dot x y| (real
  $\vx,\vy$); (d) \lstinline[style=l4style]|nrm2 x|; (e)
  \lstinline[style=l4style]|axpby| with one real and one complex scalar.

  \item \textbf{Exactness.} Explain why the promotion $\alpha \mapsto \alpha +
  0i$ introduces \emph{no} rounding error, and why this exactness is what lets
  the operator's algebraic laws be stated once rather than twice. Contrast with
  a hypothetical \texttt{float32} $\sqsubseteq$ \texttt{float64} promotion: is
  \emph{that} one exact in both directions? Which direction, if either, is safe?
\end{enumerate}
